% Stage 8 restored fragment — 2.8.2 Orthogonality and coupling of radiative modes (archive lines 10597-10908)
\subsection{Orthogonality and coupling of radiative modes}
\label{sec:ch282}

A standard-gain horn on a compact-range turntable fixes the laboratory scene for `Orthogonality and coupling of radiative modes'': frequency, probe standoff, and wattmeter closure on $\Gamma_R$ precede the formal displays below \cite{hansen1988}. Maxwell's retarded export to $\Gamma_R$ remains the governing constraint \cite{maxwell1873}.
Bistatic radar and compact-range scattering matrices rely on reciprocity to relate monostatic
calibrations to full \(S\)-parameter sets \cite{stratton1941}. The adjoint theory extends
one-port antenna reciprocity to continuous source and trace pairings \cite{devaney2004}.
Plane-wave decompositions used in reflector and radar-cross-section campaigns inherit reciprocity
constraints on amplitudes before any angular filter is applied \cite{hansen1988}. When
reciprocity appears broken after probe correction, the fault usually lies in inconsistent
impressed-current normalization or a missing outgoing tail in the Green operator, not in
Maxwell theory itself \cite{hansen1988}. Modal orthogonality and coupling on
\(\operatorname{ran}\mathcal{R}_\omega\) organize how radiative modes exchange power between
ports; reciprocity fixes the bilinear pairing before angular filters act
\cite{stratton1941}.
\medskip
Subsection~\ref{sec:ch282} gives the Chapter~\ref{ch:02} definitive account of orthogonality and
coupling of radiative modes on the spectral coordinates fixed in Subsection~\ref{sec:ch261},
developing Eqs.~\eqref{eq:ch2.modal-orthogonality}
and~\eqref{eq:ch2.radiative-coupling-bilinear} \cite{stratton1941}. Radiation eigenpairs
Eq.~\eqref{eq:ch2.radiation-eigenpair-def}, modal projectors
Eq.~\eqref{eq:ch2.modal-projector-def}, source amplitudes
Eq.~\eqref{eq:ch2.source-modal-amplitude}, spectral decomposition of the radiation operator
(Subsection~\ref{sec:ch261}), and biorthogonality
Eq.~\eqref{eq:ch2.radiation-biorthogonality} are presupposed; eigenpair construction and
reciprocity-based solvability are not re-proved \cite{stratton1941}. Explicit vector
spherical harmonic orthogonality belongs to \ref{ch:03}
Subsection~3.4.3; continuous-spectrum normalization belongs to Subsection~3.5.2, both with USE of
the present operator targets \cite{stratton1941}. Transmit--receive systems are deferred to
Subsection~\ref{sec:ch283} \cite{harrington2001}.
Assumptions. Fix \(\omega>0\), outgoing fields in \(\mathcal{F}_{\mathrm{SF}}\), radiative images in
\(\mathcal{F}_{\mathrm{rad}}\), and orthonormal radiation eigenmodes
\(\{\mathbf{u}_n\}_{n\in\mathcal{N}_{\mathrm{rad}}}\) from spectral decomposition of the radiation operator (Subsection~\ref{sec:ch261}) on the declared domain class
\cite{stratton1941}. Media satisfy reciprocal class
Eq.~\eqref{eq:ch2.reciprocal-medium-class} and passive boundary
Eq.~\eqref{eq:ch2.passive-impedance-class} so Lorentz reciprocity for the solution map (Subsection~\ref{sec:ch281}) applies \cite{stratton1941}. Adjoint modes
\(\widetilde{\mathbf{u}}_n\) are defined through
Eq.~\eqref{eq:ch2.radiation-biorthogonality} and coincide with \(\mathbf{u}_n\) on
\(\mathcal{F}_{\mathrm{rad}}\) under reciprocity-qualified pairings below
\cite{stratton1941}. regime qualification Eq.~\eqref{eq:ch2.regime-operator-screen}
governs transport interpretation of coupling reports \cite{harrington2001}.
On the orthonormal family from spectral decomposition of the radiation operator
(Subsection~\ref{sec:ch261}), distinct exporting modes are declared orthogonal with respect to the
energy inner product $\langle\cdot,\cdot\rangle_{\mathcal{U}}$ on $\mathcal{F}_{\mathrm{rad}}$.
The locked orthogonality statement reads

\begin{equation}
\label{eq:ch2.radiation-modal-orthogonality}
\langle \mathbf{u}_m,\,\mathbf{u}_n\rangle_{\mathcal{U}}
=
\delta_{mn},
\qquad
m,n\in\mathcal{N}_{\mathrm{rad}},
\end{equation}
with \(\|\mathbf{u}_n\|_{\mathcal{U}}=1\) already enforced in
Eq.~\eqref{eq:ch2.radiation-eigenpair-def} \cite{stratton1941}.
Equation~\eqref{eq:ch2.radiation-modal-orthogonality} is the Subsection~\ref{sec:ch282} locked definition of
upgrade of the first line of Eq.~\eqref{eq:ch2.modal-orthogonality}; it
specializes projector orthogonality \(P_m P_n=\delta_{mn}P_n\) from
Eq.~\eqref{eq:ch2.modal-projector-def} to the chosen orthonormal basis
\cite{stratton1941}. Algebraically, Eq.~\eqref{eq:ch2.radiation-modal-orthogonality} makes
\(\{\mathbf{u}_n\}\) an orthonormal frame on \(\mathcal{F}_{\mathrm{rad}}\) for energy inner
product Eq.~\eqref{eq:ch2.field-energy-norm}; physically, distinct exporting modes carry
non-overlapping far-zone energy content on \(S^{2}\) when
\(\mathcal{D}_{J,\mathrm{reg}}=1\) \cite{jackson1999}. Explicit angular
integrals producing Eq.~\eqref{eq:ch2.radiation-modal-orthogonality} for harmonic bases are
not tabulated here \cite{stratton1941}.
For canonical eigen-sources \(\Jimp^{(n)}\) in Eq.~\eqref{eq:ch2.radiation-eigenpair-def},
define the \emph{modal coupling coefficient}
\begin{equation}
\label{eq:ch2.modal-coupling-coefficient}
C_{mn}
=
\langle \widetilde{\mathbf{u}}_m,\,\mathcal{R}_\omega[\Jimp^{(n)}]\rangle_{\mathcal{U}},
\qquad
m,n\in\mathcal{N}_{\mathrm{rad}},
\end{equation}
with adjoint modes from Eq.~\eqref{eq:ch2.radiation-biorthogonality}
\cite{stratton1941}. Equation~\eqref{eq:ch2.modal-coupling-coefficient} is the
Subsection~\ref{sec:ch282} development of the second line of Eq.
Eq.~\eqref{eq:ch2.modal-orthogonality}. When
Eq.~\eqref{eq:ch2.radiation-biorthogonality} holds,
\begin{equation}
\label{eq:ch2.modal-coupling-diagonal}
C_{mn}=\lambda_n\,\delta_{mn}
\end{equation}
for exporting eigenpairs with \(\lambda_n\neq 0\); null modes
\(n\in\mathcal{N}_{\mathrm{null}}\) Eq.~\eqref{eq:ch2.modal-null-index-set} contribute
\(C_{mn}=0\) for all \(m\) \cite{stratton1941}. Algebraically,
Eq.~\eqref{eq:ch2.modal-coupling-diagonal} diagonalizes coupling in an eigenmode basis;
physically, orthogonal radiation modes do not co-excite one another through
\(\mathcal{R}_\omega\) when sources align with eigenpairs---cross-talk appears only from
superpositions or degenerate subspaces \cite{harrington2001}.
For arbitrary admissible \(\Jimp^{(1)},\Jimp^{(2)}\in\mathcal{J}_\Omega\), define
\begin{equation}
\label{eq:ch2.radiative-coupling-bilinear}
\mathfrak{C}_\omega\!\big(\Jimp^{(1)},\Jimp^{(2)}\big)
=
\langle \mathcal{R}_\omega[\Jimp^{(1)}],\,\mathcal{R}_\omega[\Jimp^{(2)}]\rangle_{\mathcal{U}}
\quad\text{on}\ \mathcal{F}_{\mathrm{rad}},
\end{equation}
developing Eq.~\eqref{eq:ch2.radiative-coupling-bilinear}
\cite{stratton1941}. Equation~\eqref{eq:ch2.radiative-coupling-bilinear}
measures export overlap on \(\operatorname{ran}\mathcal{R}_\omega\), not raw reactive overlap
of \(\mathcal{S}_\omega[\Jimp^{(k)}]\) \cite{stratton1941}. Under Lorentz reciprocity for the solution map (Subsection~\ref{sec:ch281})(iv),
\begin{equation}
\label{eq:ch2.coupling-bilinear-symmetry}
\mathfrak{C}_\omega\!\big(\Jimp^{(1)},\Jimp^{(2)}\big)
=
\mathfrak{C}_\omega\!\big(\Jimp^{(2)},\Jimp^{(1)}\big)
\end{equation}
on reciprocal media \cite{stratton1941}. Reactive near-zone terms in
\(\Pi_{\mathrm{ev}}\mathcal{S}_\omega[\Jimp]\) do not enter
Eq.~\eqref{eq:ch2.radiative-coupling-bilinear} through
Eq.~\eqref{eq:ch2.rad-null-on-evanescent} \cite{stratton1941}.
Let \(c_n[\Jimp^{(k)}]\) denote source amplitudes Eq.~\eqref{eq:ch2.source-modal-amplitude}
for \(k=1,2\). When the modal sum converges in \(\mathcal{U}\),
\begin{equation}
\label{eq:ch2.modal-coupling-sum}
\mathfrak{C}_\omega\!\big(\Jimp^{(1)},\Jimp^{(2)}\big)
=
\sum_{m,n\in\mathcal{N}_{\mathrm{rad}}}
\lambda_m\lambda_n\,
c_m[\Jimp^{(1)}]^{*}\,
c_n[\Jimp^{(2)}]\,
\langle \mathbf{u}_m,\mathbf{u}_n\rangle_{\mathcal{U}},
\end{equation}
using Eq.~\eqref{eq:ch2.radiation-map-modal-sum} and linearity
Eq.~\eqref{eq:ch2.radiation-linearity} \cite{stratton1941}. With
Eq.~\eqref{eq:ch2.radiation-modal-orthogonality},
\begin{equation}
\label{eq:ch2.modal-coupling-sum-diagonal}
\mathfrak{C}_\omega\!\big(\Jimp^{(1)},\Jimp^{(2)}\big)
=
\sum_{n\in\mathcal{N}_{\mathrm{rad}}}
|\lambda_n|^{2}\,
c_n[\Jimp^{(1)}]^{*}\,
c_n[\Jimp^{(2)}],
\end{equation}
for real nonnegative exporting gains \(\lambda_n>0\) \cite{stratton1941}. Equation
\eqref{eq:ch2.modal-coupling-sum-diagonal} shows that coupling is a weighted inner product
of modal excitation vectors on the exporting sector; geometric constraints
geometric constraints on radiative spectra (Subsection~\ref{sec:ch264}) may force individual \(c_n=0\) without
breaking Eq.~\eqref{eq:ch2.coupling-bilinear-symmetry} \cite{harrington2001}.
\begin{figure}[t]
\centering
\ChOneFigBlock{%
\begin{tikzpicture}[ch1fig, node distance=5mm and 7mm]
 \node[ch1box] (j1) {$\Jimp^{(1)}$};
 \node[ch1box, right=7mm of j1] (r) {$\mathcal{R}_\omega$};
 \node[ch1box, right=7mm of r] (u) {$\sum_n \lambda_n c_n \mathbf{u}_n$};
 \node[ch1box, below=5mm of j1] (j2) {$\Jimp^{(2)}$};
 \node[ch1box, below=5mm of u] (c) {$\mathfrak{C}_\omega$};
 \node[ch1box, left=7mm of c] (orth) {Eq.~\eqref{eq:ch2.radiation-modal-orthogonality}};
 \draw[ch1flow] (j1.east) -- (j1.east |- r.west);
 \draw[ch1flow] (j1.east |- r.west) -- (r.west);
 \draw[ch1flow] (j2.east) -- (j2.east |- r.west);
 \draw[ch1flow] (j2.east |- r.west) -- (r.west);
 \draw[ch1flow] (r.east) -- (r.east |- u.west);
 \draw[ch1flow] (r.east |- u.west) -- (u.west);
 \draw[ch1flow] (u.south) -- (u.south |- c.north);
 \draw[ch1flow] (u.south |- c.north) -- (c.north);
 \draw[ch1flow] (orth.west) -- (orth.west |- c.east);
 \draw[ch1flow] (orth.west |- c.east) -- (c.east);
\end{tikzpicture}%
}
\caption{Subsection~\ref{sec:ch282} modal coupling measurement sequence: sources map through
\(\mathcal{R}_\omega\) to orthonormal modes Eq.~\eqref{eq:ch2.radiation-modal-orthogonality};
bilinear coupling Eq.~\eqref{eq:ch2.radiative-coupling-bilinear} factorizes through modal
amplitudes Eq.~\eqref{eq:ch2.modal-coupling-sum}.}
\label{fig:ch2.modal-coupling-measurement chain}
\end{figure}
Figure interpretation. Two impressed sources $\Jimp^{(1)}$, $\Jimp^{(2)}$ share the same radiation map $\mathcal{R}_\omega$ and expand on orthonormal modes Eq.~\eqref{eq:ch2.radiation-modal-orthogonality}. The coupling node $\mathfrak{C}_\omega$ measures export overlap on $\mathcal{F}_{\mathrm{rad}}$, not raw near-zone reactive cross terms. Orthogonality diagonalizes coupling when sources align with eigenpairs; cross-talk appears only from superposition or degenerate subspaces. Modes with $\lambda_n=0$ feed the null sector without contributing to $\mathfrak{C}_\omega$ on the range \cite{stratton1941}. Figure~\ref{fig:ch2.modal-coupling-measurement chain} separates orthogonality of mode shapes from
bilinear coupling of impressed sources: Eq.~\eqref{eq:ch2.radiation-modal-orthogonality} is
intrinsic to \(\mathcal{F}_{\mathrm{rad}}\); Eq.~\eqref{eq:ch2.radiative-coupling-bilinear}
depends on both sources \cite{stratton1941}. On a reciprocal two-mode truncation with $\lambda_1=0.42$, $\lambda_2=0.31$ and orthonormal $\mathbf{u}_1$, $\mathbf{u}_2$, excite $\Jimp^{(1)}=c_1\mathbf{J}^{(1)}_{\mathrm{imp}}$ and probe with $\Jimp^{(2)}=c_2\mathbf{J}^{(2)}_{\mathrm{imp}}$. Eq.~\eqref{eq:ch2.modal-coupling-sum-diagonal} gives $\mathfrak{C}_\omega(\Jimp^{(1)},\Jimp^{(2)})=\lambda_1\lambda_2 c_1^* c_2$. Swapping labels leaves the scalar unchanged, so normalized $S_{21}=S_{12}$ on the exporting sector. Adding $\Jimp^{\mathrm{(NR)}}$ to either port shifts $\mathcal{S}_\omega[\Jimp]$ near $\Omega_s$ but leaves $S_{21}$ fixed---the walkthrough confirms reciprocity on $\operatorname{ran}\mathcal{R}_\omega$, not on raw near-field coupling matrices \cite{stratton1941}.
\phantomsection\label{thm:ch2.radiation-modal-orthogonality-coupling}
Let \(\{\mathbf{u}_n,\Jimp^{(n)},\lambda_n\}\) be radiation eigenpairs
Eq.~\eqref{eq:ch2.radiation-eigenpair-def} forming an orthonormal basis of
\(\mathcal{F}_{\mathrm{rad}}\) on the declared domain class, with reciprocal media
Eq.~\eqref{eq:ch2.reciprocal-medium-class}. Then:
\begin{enumerate}
\item[(i)] Orthogonality Eq.~\eqref{eq:ch2.radiation-modal-orthogonality} holds on
\(\mathcal{N}_{\mathrm{rad}}\).
\item[(ii)] Coupling coefficients satisfy Eq.~\eqref{eq:ch2.modal-coupling-diagonal} and
biorthogonality Eq.~\eqref{eq:ch2.radiation-biorthogonality}.
\item[(iii)] For all admissible \(\Jimp^{(1)},\Jimp^{(2)}\),
Eqs.~\eqref{eq:ch2.radiative-coupling-bilinear}--\eqref{eq:ch2.modal-coupling-sum-diagonal}
hold, with symmetry Eq.~\eqref{eq:ch2.coupling-bilinear-symmetry} from Lorentz reciprocity for the solution map (Subsection~\ref{sec:ch281}).
\item[(iv)] Adding \(\Jimp^{\mathrm{(NR)}}\in\ker\mathcal{R}_\omega\) to either source leaves
Eq.~\eqref{eq:ch2.radiative-coupling-bilinear} unchanged.
\end{enumerate}
Part (i) is normalization and completeness from spectral decomposition of the radiation operator (Subsection~\ref{sec:ch261}) and
Eq.~\eqref{eq:ch2.radiation-spectral-expansion} \cite{stratton1941}. Part (ii) substitutes
\(\mathcal{R}_\omega[\Jimp^{(n)}]=\lambda_n\mathbf{u}_n\) into
Eq.~\eqref{eq:ch2.modal-coupling-coefficient} and applies
Eq.~\eqref{eq:ch2.radiation-biorthogonality} \cite{stratton1941}. Part (iii)
expands \(\mathcal{R}_\omega[\Jimp^{(k)}]\) via
Eq.~\eqref{eq:ch2.radiation-map-modal-sum}; symmetry is Lorentz reciprocity for the solution map (Subsection~\ref{sec:ch281})(iv) \cite{stratton1941}. Part (iv) follows from
\(\mathcal{R}_\omega[\Jimp^{\mathrm{(NR)}}]=\mathbf{0}\) and
Eq.~\eqref{eq:ch2.nr-superposition-law} \cite{stratton1941}. The reciprocal reduction of adjoint modes proceeds as follows: under Eqs.~\eqref{eq:ch2.reciprocal-medium-class}--\eqref{eq:ch2.passive-impedance-class},
\begin{equation}
\label{eq:ch2.reciprocal-adjoint-mode-identification}
\widetilde{\mathbf{u}}_n=\mathbf{u}_n
\quad\text{on}\ \mathcal{F}_{\mathrm{rad}}
\quad\text{for exporting modes with}\ \lambda_n\neq 0,
\end{equation}
so Eq.~\eqref{eq:ch2.radiation-biorthogonality} reduces to standard orthogonality
Eq.~\eqref{eq:ch2.radiation-modal-orthogonality} with diagonal coupling
Eq.~\eqref{eq:ch2.modal-coupling-diagonal} \cite{stratton1941}. Equation
\eqref{eq:ch2.reciprocal-adjoint-mode-identification} is the operator form of reciprocity for
modal probes: receive patterns aligned with \(\mathbf{u}_n\) coincide with transmit modes
when media and boundaries are reciprocal \cite{harrington2001}. Non-reciprocal
media break Eq.~\eqref{eq:ch2.reciprocal-adjoint-mode-identification}; biorthogonality with
distinct \(\widetilde{\mathbf{u}}_n\) remains the general template imported from
Subsection~\ref{sec:ch261} \cite{stratton1941}. When \(\lambda_n=\lambda_{n'}\) for \(n\neq n'\) within a rotationally degenerate family from
Subsection~\ref{sec:ch262}, Eq.~\eqref{eq:ch2.modal-coupling-diagonal} still holds for
\emph{eigen} sources \(\Jimp^{(n)}\), but superpositions within the degenerate block carry
off-diagonal coupling in an arbitrary basis. Let \(\mathcal{I}_\ell\subset\mathcal{N}_{\mathrm{rad}}\)
index a \((2\ell+1)\)-fold degenerate family. Any orthonormal basis
\(\{\mathbf{u}_n\}_{n\in\mathcal{I}_\ell}\) spanning the block satisfies
Eq.~\eqref{eq:ch2.radiation-modal-orthogonality}, yet a fixed physical source may excite
multiple members with coefficients tied to orientation through
Eq.~\eqref{eq:ch2.radiation-rotation-commutation} \cite{stratton1941}. Define
the block coupling matrix
\begin{equation}
\label{eq:ch2.degenerate-block-coupling-matrix}
\mathbf{K}^{(\ell)}_{mn}
=
\langle \mathbf{u}_m,\,\mathcal{R}_\omega[\Jimp^{(n)}]\rangle_{\mathcal{U}},
\qquad
m,n\in\mathcal{I}_\ell,
\end{equation}
which is diagonal in the eigenvector basis of \(\Jimp^{(n)}\) but not in laboratory-fixed
coordinates \cite{stratton1941}. Reciprocity
Eq.~\eqref{eq:ch2.coupling-bilinear-symmetry} forces \(\mathbf{K}^{(\ell)}\) to be symmetric
under Lorentz reciprocity for the solution map (Subsection~\ref{sec:ch281}) \cite{stratton1941}.
Explicit Wigner coupling of block entries belongs to \ref{ch:03}
\cite{stratton1941}.
\phantomsection\label{prop:ch2.superposition-cross-coupling}
If \(\Jimp^{(1)}=\sum_{n\in\mathcal{N}_{\mathrm{rad}}}a_n\Jimp^{(n)}\) and
\(\Jimp^{(2)}=\sum_{m\in\mathcal{N}_{\mathrm{rad}}}b_m\Jimp^{(m)}\) are admissible finite
superpositions, then
\begin{equation}
\label{eq:ch2.superposition-coupling-expansion}
\mathfrak{C}_\omega\!\big(\Jimp^{(1)},\Jimp^{(2)}\big)
=
\sum_{m,n\in\mathcal{N}_{\mathrm{rad}}}
a_m^{*}b_n\,\lambda_m\lambda_n\,
\langle \mathbf{u}_m,\mathbf{u}_n\rangle_{\mathcal{U}},
\end{equation}
reducing to Eq.~\eqref{eq:ch2.modal-coupling-sum-diagonal} when \(\langle\mathbf{u}_m,\mathbf{u}_n\rangle=\delta_{mn}\)
\cite{stratton1941}. Off-diagonal terms \(m\neq n\) appear only when at least
one source superposes non-collinear modes; single-mode excitation yields diagonal coupling
\cite{jackson1999}.
Linearity Eq.~\eqref{eq:ch2.radiation-linearity} and
Eq.~\eqref{eq:ch2.radiation-map-modal-sum} \cite{stratton1941}.
\phantomsection\label{prop:ch2.observation-limited-coupling}
For observation map \(\mathcal{M}_\omega=\mathcal{O}_{\mathrm{meas}}\circ\Lambda_{\mathrm{rad}}\)
Eq.~\eqref{eq:ch2.reconstruction-observation-map}, define observed coupling
\(\mathfrak{C}_{\omega,\mathrm{obs}}=
\langle \mathcal{O}_{\mathrm{meas}}[\mathcal{R}_\omega[\Jimp^{(1)}]],\,
\mathcal{O}_{\mathrm{meas}}[\mathcal{R}_\omega[\Jimp^{(2)}]]\rangle\).
Then \(\mathfrak{C}_{\omega,\mathrm{obs}}\) is symmetric under Lorentz reciprocity for the solution map (Subsection~\ref{sec:ch281}) when \(\mathcal{O}_{\mathrm{meas}}\) is self-adjoint
on \(\mathcal{F}_{\mathrm{rad}}\), but may vanish for pairs with
\(\mathcal{R}_\omega[\Jimp^{(k)}]\in\ker\mathcal{O}_{\mathrm{meas}}\) even when
Eq.~\eqref{eq:ch2.radiative-coupling-bilinear} is nonzero \cite{stratton1941}.
Rank cap Eq.~\eqref{eq:ch2.reconstruction-rank-cap} limits the number of independently
observable coupling channels \cite{stratton1941}.
Apply Eq.~\eqref{eq:ch2.radiation-map-reciprocity} and self-adjointness of
\(\mathcal{O}_{\mathrm{meas}}\) \cite{stratton1941}. Kernel inclusion
\(\ker\mathcal{R}_\omega\subset\ker\mathcal{M}_\omega\) is
measurement null-space extension (Subsection~\ref{sec:ch271}) \cite{stratton1941}.
\begin{equation}
\label{eq:ch2.kernel-field-image-law}
\Jimp\in\ker\mathcal{R}_\omega
\ \Longrightarrow\
\mathcal{S}_\omega[\Jimp]\in\mathcal{F}_{\mathrm{SF}},
\qquad
\Pi_{\mathrm{rad}}\,\mathcal{S}_\omega[\Jimp]=\mathbf{0},
\end{equation}
Accessible modal sector Eq.~\eqref{eq:ch2.accessible-modal-sector} and angular budget
Eq.~\eqref{eq:ch2.angular-spectrum-composition-law} restrict which terms in
Eq.~\eqref{eq:ch2.modal-coupling-sum-diagonal} are observable in far-zone reports
\cite{stratton1941}. Null modes \(\lambda_n=0\) carry zero coupling diagonal
entries but may alter near reactive content through \(\mathcal{S}_\omega\)
Eq.~\eqref{eq:ch2.kernel-field-image-law} without changing
Eq.~\eqref{eq:ch2.radiative-coupling-bilinear} \cite{stratton1941}. Composition with
null-space invariance of radiative reciprocity (Subsection~\ref{sec:ch281}) and
orthogonality and coupling of radiative modes (Subsection~\ref{sec:ch282})(iv) explains why invisible
currents do not break reciprocity-qualified coupling on \(\operatorname{ran}\mathcal{R}_\omega\)
\cite{harrington2001}.
Orthogonality and coupling of radiative modes (Subsection~\ref{sec:ch282}) requires linear reciprocal media
at fixed \(\omega\), outgoing radiation condition, domain-class completeness of
\(\{\mathbf{u}_n\}\), and convergent modal sums \cite{stratton1941}. Loss
moves \(\lambda_n\) into the complex plane; Eq.~\eqref{eq:ch2.modal-coupling-sum-diagonal}
generalizes with complex conjugation on one amplitude slot \cite{stratton1941}. Truncated
modal books on cavities do not substitute for exterior continuous spectra
Eq.~\eqref{eq:ch2.exterior-continuous-spectral-integral} when coupling far-zone channels
\cite{jackson1999}. regime qualification Eq.~\eqref{eq:ch2.regime-component-criterion} is
mandatory before interpreting Eq.~\eqref{eq:ch2.radiative-coupling-bilinear} as transport
\cite{stratton1941}.
Subsections~\ref{sec:ch281}--\ref{sec:ch282} established Lorentz reciprocity, radiative bilinear
coupling Eq.~\eqref{eq:ch2.radiative-coupling-bilinear}, and modal orthogonality on
\(\mathcal{F}_{\mathrm{rad}}\) \cite{stratton1941}. The final reciprocity link ties those
operator statements to finite ports: transmit and receive maps on reciprocal hardware must
interchange when excitation and probe roles are swapped.
